% !TEX root = ../main.tex
\begin{ouverture}[redox]{3}{Réactions d'oxydoréduction}{Partie 1 · Constitution et transformations de la matière}
{Oxydant\enspace·\enspace Réducteur\enspace·\enspace Couple\enspace·\enspace Demi-équation\enspace·\enspace Transfert d'électrons}
\objectif{Identifier, à partir de données expérimentales, un transfert d'électrons entre deux réactifs.}
\objectif{Définir oxydant, réducteur, oxydation, réduction et couple oxydant-réducteur.}
\objectif{Écrire la demi-équation électronique d'un couple, y compris en milieu acide.}
\objectif{Établir l'équation d'une réaction d'oxydoréduction, les couples étant donnés.}
\objectif{Vérifier la conservation des éléments et des charges dans une équation, y compris avec Python.}
\prerequis{Ions monoatomiques et polyatomiques, formules chimiques (classe de seconde).}
\prerequis{Modélisation d'une transformation par une équation de réaction ajustée (classe de seconde).}
\prerequis{Quantité de matière et volume molaire (chapitre~1).}
\end{ouverture}
\essentiel
\essSection{1}{Des transformations par transfert d'électrons}
\begin{definition}[Oxydant, réducteur, couple]
Un \textbf{oxydant} est une espèce capable de capter un ou plusieurs électrons ; un \textbf{réducteur}, une espèce capable d'en
céder. Un oxydant et un réducteur qui se transforment l'un en l'autre par transfert d'électrons forment un \textbf{couple
oxydant-réducteur}, noté Ox/Red. On le décrit par une \textbf{demi-équation électronique} :
\[\mathrm{Ox}+n\,\mathrm{e}^-=\mathrm{Red}\qquad\text{par exemple}\qquad\ce{Cu^{2+} + 2e- = Cu}.\]
Le signe $=$ indique que l'échange peut se faire dans les deux sens.
\end{definition}
\schema{03-zinc}
\begin{definition}[Oxydation et réduction]
Une \textbf{oxydation} est une perte d'électrons, une \textbf{réduction} un gain d'électrons. Au cours d'une réaction
d'oxydoréduction, le réducteur d'un couple est oxydé, et l'oxydant d'un autre couple est réduit : les électrons cédés
par l'un sont captés par l'autre, et n'apparaissent jamais libres en solution.
\end{definition}
\essSection{2}{Écrire une demi-équation électronique}
\begin{savoirfaire}[Ajuster une demi-équation en milieu acide]
Exemple du couple $\ce{MnO4-}/\ce{Mn^{2+}}$ :
\begin{questions}
\item Écrire l'oxydant et le réducteur de part et d'autre du signe $=$ et ajuster les éléments autres que O et H :
$\ce{MnO4- = Mn^{2+}}$.
\item Ajuster les atomes d'oxygène avec des molécules d'eau : $\ce{MnO4- = Mn^{2+} + 4H2O}$.
\item Ajuster les atomes d'hydrogène avec des ions $\ce{H+}$ : $\ce{MnO4- + 8H+ = Mn^{2+} + 4H2O}$.
\item Ajuster les charges avec des électrons, du côté de l'oxydant : $\ce{MnO4- + 8H+ + 5e- = Mn^{2+} + 4H2O}$.
\end{questions}
\end{savoirfaire}
\essSection{3}{Établir l'équation d'une réaction d'oxydoréduction}
\begin{propriete}[Combiner deux demi-équations]
On écrit la demi-équation de l'oxydant qui réagit, dans le sens de sa réduction, et celle du réducteur, dans le sens de
son oxydation. On multiplie l'une ou l'autre pour que le nombre d'électrons cédés soit égal au nombre d'électrons captés,
puis on additionne : les électrons disparaissent de l'équation.
\end{propriete}
\begin{exemple}
Les ions permanganate oxydent les ions fer(II) : $\ce{MnO4- + 8H+ + 5e- -> Mn^{2+} + 4H2O}$ et $\ce{Fe^{2+} -> Fe^{3+} + e-}$, à
multiplier par 5. En additionnant :
\[\ce{MnO4- + 8H+ + 5Fe^{2+} -> Mn^{2+} + 5Fe^{3+} + 4H2O}.\]
On vérifie la conservation des éléments (1 Mn, 4 O, 8 H, 5 Fe) et des charges ($+17$ de chaque côté).
\end{exemple}
\begin{python}
especes = {'MnO4-': ({'Mn': 1, 'O': 4}, -1), 'H+': ({'H': 1}, 1),
           'Fe2+': ({'Fe': 1}, 2), 'Mn2+': ({'Mn': 1}, 2),
           'Fe3+': ({'Fe': 1}, 3), 'H2O': ({'H': 2, 'O': 1}, 0)}
reactifs = {'MnO4-': 1, 'H+': 8, 'Fe2+': 5}
produits = {'Mn2+': 1, 'Fe3+': 5, 'H2O': 4}

def bilan(cote):
    atomes, charge = {}, 0
    for nom, coef in cote.items():
        composition, z = especes[nom]
        charge = charge + coef * z
        for el, nb in composition.items():
            atomes[el] = atomes.get(el, 0) + coef * nb
    return atomes, charge

print(bilan(reactifs))
print(bilan(produits))
\end{python}
\begin{console}
({'Mn': 1, 'O': 4, 'H': 8, 'Fe': 5}, 17)
({'Mn': 1, 'Fe': 5, 'H': 8, 'O': 4}, 17)
\end{console}
\noindent\capacite{Capacité numérique} Chaque espèce est décrite par sa composition et sa charge. La fonction \code{bilan} compte les atomes de
chaque élément et la charge totale d'un côté de l'équation, en tenant compte des coefficients. Les deux lignes affichées
contiennent les mêmes nombres d'atomes et la même charge : l'équation est bien ajustée.
\essSection{4}{Identifier un transfert d'électrons}
\begin{propriete}[Indices expérimentaux]
Un transfert d'électrons se repère par l'apparition ou la disparition d'espèces : dépôt métallique, dégagement gazeux,
changement de couleur d'une solution. On identifie les produits par des tests : les ions métalliques donnent avec la soude
des précipités colorés, le diiode bleuit l'empois d'amidon, le dihydrogène produit une petite détonation au contact d'une
flamme.
\end{propriete}
\schema{03-tests}
\begin{attention}
Une demi-équation n'est pas une réaction : elle décrit un couple, et les électrons n'existent jamais seuls en solution. Les
ions $\ce{H+}$ et les molécules d'eau d'une demi-équation servent à ajuster O et H en milieu acide. Dans l'équation finale, il ne
doit rester aucun électron.
\end{attention}
\begin{experience}[Le zinc et les ions cuivre(II)]
On plonge une lame de zinc décapée dans une solution bleue de sulfate de cuivre. Après quelques minutes, la lame se
couvre d'un dépôt rouge et la solution se décolore. Un prélèvement de solution, additionné de soude, donne un précipité
blanc : des ions zinc(II) sont apparus. On en déduit que le zinc a cédé des électrons aux ions cuivre(II).
\end{experience}
\begin{rediger}
Nommer l'oxydant et le réducteur qui réagissent et leurs couples. Écrire les deux demi-équations dans le sens où elles se
produisent, les multiplier pour égaliser les électrons, additionner, puis vérifier la conservation des éléments et des
charges.
\end{rediger}
\exercices

\rubrique{Découverte}
\begin{blocLivre}{ex-3-1}
\begin{exercice}[1]{Oxydants, réducteurs et couples}
\begin{questions}
\item Dans les couples $\ce{Ag+}/\ce{Ag}$, $\ce{Fe^{3+}}/\ce{Fe^{2+}}$ et $\ce{Al^{3+}}/\ce{Al}$, identifier l'oxydant et le réducteur.
\item Écrire les demi-équations électroniques de ces trois couples.
\item Écrire les demi-équations des couples $\ce{I2}/\ce{I-}$ et $\ce{Cl2}/\ce{Cl-}$.
\item L'ion fer(II) appartient à deux couples, $\ce{Fe^{3+}}/\ce{Fe^{2+}}$ et $\ce{Fe^{2+}}/\ce{Fe}$. Quel est son rôle dans chacun ?
\end{questions}
\end{exercice}
\end{blocLivre}
\begin{blocLivre}{ex-3-2}
\begin{exercice}[1]{Des demi-équations en milieu acide}
Écrire les demi-équations électroniques des couples suivants, en milieu acide.
\begin{questions}
\item $\ce{Cr2O7^{2-}}/\ce{Cr^{3+}}$ (ion dichromate et ion chrome(III)).
\item $\ce{O2}/\ce{H2O}$ et $\ce{H2O2}/\ce{H2O}$.
\item $\ce{NO3-}/\ce{NO}$ (ion nitrate et monoxyde d'azote).
\item $\ce{O2}/\ce{H2O2}$. Quel double rôle le peroxyde d'hydrogène $\ce{H2O2}$ peut-il jouer ?
\end{questions}
\end{exercice}
\end{blocLivre}
\begin{blocLivre}{ex-3-3}
\begin{exercice}[1]{L'arbre de Diane}
On plonge un fil de cuivre dans une solution incolore de nitrate d'argent, $\ce{Ag+}+\ce{NO3-}$.
\schema{03-diane}
\begin{questions}
\item Quelles observations montrent qu'une transformation chimique a lieu ?
\item Identifier les deux couples mis en jeu et écrire leurs demi-équations.
\item Établir l'équation de la réaction. Quelle espèce est oxydée ? Laquelle est réduite ?
\item Pourquoi la solution bleuit-elle ? Proposer un test pour le confirmer.
\end{questions}
\end{exercice}
\end{blocLivre}
\begin{blocLivre}{ex-3-4}
\begin{exercice}[1]{Un métal attaqué par un acide}
On verse de l'acide chlorhydrique, $\ce{H+}+\ce{Cl-}$, en excès sur \qty{0,654}{g} de zinc en poudre : un gaz se dégage, et le zinc
disparaît.
\begin{donnees}
$M(\ce{Zn})=\qty{65,4}{g.mol^{-1}}$ ; $V_\mathrm{m}=\qty{24,0}{L.mol^{-1}}$ à \qty{20}{\degreeCelsius}. Couples : $\ce{H+}/\ce{H2}$ et $\ce{Zn^{2+}}/\ce{Zn}$.
\end{donnees}
\begin{questions}
\item Comment identifier le gaz dégagé ?
\item Établir l'équation de la réaction.
\item Calculer la quantité de matière de zinc, puis celle de dihydrogène formé.
\item Calculer le volume de dihydrogène obtenu à \qty{20}{\degreeCelsius}.
\end{questions}
\end{exercice}
\end{blocLivre}
\begin{blocLivre}{ex-3-5}
\begin{exercice}[1]{Un clou dans une solution bleue}
Un clou en fer est plongé dans une solution de sulfate de cuivre. Une heure plus tard, il est couvert d'un dépôt rouge et
la solution a verdi ; un prélèvement de solution donne, avec la soude, un précipité vert.
\begin{questions}
\item Identifier le dépôt rouge et l'ion responsable du précipité vert.
\item En déduire les deux couples mis en jeu.
\item Établir l'équation de la réaction.
\item Pourquoi la solution, d'abord bleue, a-t-elle verdi ?
\end{questions}
\end{exercice}
\end{blocLivre}
\rubrique{Entraînement}
\begin{blocLivre}{ex-3-6}
\begin{exercice}[2]{Établir des équations}
Établir l'équation de chaque réaction, en milieu acide si nécessaire.
\begin{questions}
\item Les ions permanganate $\ce{MnO4-}$ oxydent les ions fer(II) ; couples $\ce{MnO4-}/\ce{Mn^{2+}}$ et $\ce{Fe^{3+}}/\ce{Fe^{2+}}$.
\item Le diiode oxyde les ions thiosulfate $\ce{S2O3^{2-}}$ ; couples $\ce{I2}/\ce{I-}$ et $\ce{S4O6^{2-}}/\ce{S2O3^{2-}}$.
\item Le peroxyde d'hydrogène oxyde les ions iodure ; couples $\ce{H2O2}/\ce{H2O}$ et $\ce{I2}/\ce{I-}$.
\item Pour chaque réaction, vérifier la conservation des charges.
\end{questions}
\end{exercice}
\end{blocLivre}
\begin{blocLivre}{ex-3-7}
\begin{exercice}[2]{L'alcootest chimique}
Dans un ancien alcootest, l'air expiré traverse un tube contenant des cristaux orange de dichromate de potassium en milieu
acide. L'éthanol de l'haleine est oxydé en acide éthanoïque, tandis que les ions dichromate, orange, sont réduits en ions
chrome(III), verts.
\begin{donnees}
Couples : $\ce{Cr2O7^{2-}}/\ce{Cr^{3+}}$ et $\ce{CH3COOH}/\ce{CH3CH2OH}$.
\end{donnees}
\begin{questions}
\item Écrire la demi-équation du couple $\ce{Cr2O7^{2-}}/\ce{Cr^{3+}}$.
\item Écrire la demi-équation du couple $\ce{CH3COOH}/\ce{CH3CH2OH}$.
\item Établir l'équation de la réaction et vérifier qu'elle est ajustée.
\item Comment le test signale-t-il la présence d'alcool ? Pourquoi la longueur de tube devenue verte renseigne-t-elle sur la
quantité d'alcool ?
\end{questions}
\end{exercice}
\end{blocLivre}
\begin{blocLivre}{ex-3-8}
\begin{exercice}[2]{Nettoyer l'argenterie}
Le noircissement de l'argenterie est dû à une fine couche de sulfure d'argent, $\ce{Ag2S}$. Pour la nettoyer, on la pose sur
une feuille d'aluminium au fond d'une bassine d'eau chaude salée : l'argent redevient brillant.
\begin{donnees}
Couples : $\ce{Ag2S}/\ce{Ag}$, de demi-équation $\ce{Ag2S + 2e- = 2Ag + S^{2-}}$, et $\ce{Al^{3+}}/\ce{Al}$. Masses molaires (\unit{g.mol^{-1}}) :
Al 27,0 ; S 32,1 ; Ag 107,9.
\end{donnees}
\begin{questions}
\item Quel est l'oxydant ? Quel est le réducteur ?
\item Établir l'équation de la réaction.
\item Une cuillère porte \qty{0,100}{g} de sulfure d'argent. Calculer la masse d'argent régénérée.
\item Calculer la masse minimale d'aluminium consommée. Pourquoi ce procédé ne fait-il pas perdre d'argent, contrairement
à un polissage ?
\end{questions}
\end{exercice}
\end{blocLivre}
\begin{blocLivre}{ex-3-9}
\begin{exercice}[2]{Vérifier une équation en Python}
On reprend le programme de l'essentiel, qui vérifie l'équation de la réaction entre les ions permanganate et fer(II).
\begin{questions}
\item \capacite{Capacité numérique}Que contient le dictionnaire \code{especes} ? Que renvoie la fonction \code{bilan} ?
\item Un élève remplace \code{'H+': 8} par \code{'H+': 4} dans \code{reactifs}. Que devient la première ligne affichée ? Que peut-on en
conclure ?
\item Quelles lignes faut-il modifier pour vérifier l'équation $\ce{Zn + Cu^{2+} -> Zn^{2+} + Cu}$ ? Écrire les dictionnaires obtenus.
\item Pourquoi faut-il comparer à la fois les atomes et les charges ?
\end{questions}
\end{exercice}
\end{blocLivre}
\begin{blocLivre}{ex-3-10}
\begin{exercice}[2]{La vitamine C, un antioxydant}
La vitamine C, ou acide ascorbique $\ce{C6H8O6}$, est le réducteur du couple $\ce{C6H6O6}/\ce{C6H8O6}$. On l'ajoute à certains aliments
comme antioxydant (E300).
\begin{questions}
\item Écrire la demi-équation du couple de l'acide ascorbique.
\item L'acide ascorbique réduit le diiode : établir l'équation de la réaction, avec le couple $\ce{I2}/\ce{I-}$.
\item Il réagit aussi avec le dioxygène de l'air, couple $\ce{O2}/\ce{H2O}$. Établir l'équation de cette réaction.
\item Expliquer pourquoi on parle d'antioxydant, et pourquoi un jus d'orange perd sa vitamine C s'il reste ouvert à l'air.
\end{questions}
\end{exercice}
\end{blocLivre}
\begin{blocLivre}{ex-3-11}
\begin{exercice}[2]{Ne jamais mélanger eau de Javel et détartrant}
L'eau de Javel contient des ions hypochlorite $\ce{ClO-}$ et des ions chlorure $\ce{Cl-}$. Mélangée à un détartrant acide, elle libère du
dichlore, un gaz toxique.
\begin{donnees}
Couples : $\ce{ClO-}/\ce{Cl2}$ et $\ce{Cl2}/\ce{Cl-}$. Une eau de Javel à 2,6~\% de chlore actif contient environ \qty{0,37}{mol.L^{-1}} d'ions hypochlorite.
$V_\mathrm{m}=\qty{24,0}{L.mol^{-1}}$ ; $M(\ce{Cl})=\qty{35,5}{g.mol^{-1}}$.
\end{donnees}
\begin{questions}
\item Écrire les demi-équations des deux couples, en milieu acide.
\item Montrer que l'équation de la réaction s'écrit $\ce{ClO- + Cl- + 2H+ -> Cl2 + H2O}$. Quelle particularité présente-t-elle ?
\item On verse par erreur \qty{100}{mL} d'eau de Javel dans un détartrant acide en excès. Calculer la quantité, le volume à \qty{20}{\degreeCelsius} et la
masse de dichlore qui peuvent se former.
\item Pourquoi le pictogramme « Ne pas mélanger avec d'autres produits » figure-t-il sur les flacons ?
\end{questions}
\end{exercice}
\end{blocLivre}
\rubrique{Objectif DS}
\begin{blocLivre}{ex-3-12}
\begin{exercice}[3]{Protéger une coque de bateau}
\begin{docu}{1}{La corrosion}
Au contact de l'eau de mer, l'acier d'une coque, principalement constitué de fer, est oxydé par le dioxygène dissous :
le fer se transforme en ions fer(II), puis en rouille.
\end{docu}
\begin{docu}{2}{Les anodes sacrificielles}
On fixe sur la coque des blocs de zinc, en contact électrique avec l'acier. Le zinc, plus réducteur que le fer, est oxydé à
sa place : il se dissout lentement en ions zinc(II), et le fer est protégé. Un bateau consomme ainsi environ \qty{2,0}{kg} de
zinc par an.
\end{docu}
\begin{donnees}
Couples : $\ce{Fe^{2+}}/\ce{Fe}$, $\ce{Zn^{2+}}/\ce{Zn}$, $\ce{O2}/\ce{H2O}$ ; on raisonne en milieu acide. Masses molaires : Fe 55,8 ; Zn 65,4 (\unit{g.mol^{-1}}).
\end{donnees}
\begin{questions}
\item Établir l'équation de l'oxydation du fer par le dioxygène.
\item Établir l'équation de l'oxydation du zinc par le dioxygène.
\item Calculer la quantité de zinc consommée en un an, puis la quantité d'électrons qu'il cède.
\item \initiative Quelle masse de fer la coque aurait-elle perdue en un an sans les anodes, si les mêmes électrons avaient
été cédés par le fer ?
\end{questions}
\end{exercice}
\end{blocLivre}
\begin{blocLivre}{ex-3-13}
\begin{exercice}[3]{L'eau oxygénée « à 10 volumes »}
\begin{docu}{1}{Une eau oxygénée qui mousse}
Versée sur une plaie, l'eau oxygénée (solution de peroxyde d'hydrogène $\ce{H2O2}$) mousse : une enzyme du sang accélère sa
transformation en eau et en dioxygène. Le peroxyde d'hydrogène appartient aux couples $\ce{H2O2}/\ce{H2O}$ et $\ce{O2}/\ce{H2O2}$.
\end{docu}
\begin{docu}{2}{Le titre en volumes}
Une eau oxygénée à 10 volumes libère, en se décomposant entièrement, \qty{10}{L} de dioxygène par litre de solution,
volume mesuré à \qty{0}{\degreeCelsius} et \qty{1013}{hPa}, conditions où le volume molaire vaut \qty{22,4}{L.mol^{-1}}.
\end{docu}
\begin{questions}
\item Écrire les demi-équations des deux couples du peroxyde d'hydrogène.
\item Montrer que l'équation de la décomposition s'écrit $\ce{2H2O2 -> 2H2O + O2}$. Quel double rôle joue le peroxyde d'hydrogène ?
\item \initiative Montrer que la concentration en peroxyde d'hydrogène d'une eau oxygénée à 10 volumes vaut environ
\qty{0,89}{mol.L^{-1}}.
\item Quel volume de dioxygène un flacon de \qty{250}{mL} peut-il libérer à \qty{0}{\degreeCelsius}, puis à \qty{20}{\degreeCelsius} ($V_\mathrm{m}=\qty{24,0}{L.mol^{-1}}$) ?
\end{questions}
\end{exercice}
\end{blocLivre}
\begin{blocLivre}{ex-3-14}
\begin{exercice}[3]{Récupérer l'argent d'un bain photographique}
\begin{docu}{1}{Un déchet précieux}
Les bains de fixateur des laboratoires photographiques contiennent des ions argent. On les récupère en plongeant dans le
bain de la poudre ou de la laine de zinc : l'argent se dépose, et le zinc passe en solution.
\end{docu}
\begin{docu}{2}{Le bain}
Un bain usagé de \qty{5,0}{L} contient \qty{2,0}{g} d'argent par litre, sous forme d'ions $\ce{Ag+}$.
\end{docu}
\begin{donnees}
Couples : $\ce{Ag+}/\ce{Ag}$ et $\ce{Zn^{2+}}/\ce{Zn}$ ; Ag 107,9 ; Zn 65,4 (\unit{g.mol^{-1}}).
\end{donnees}
\begin{questions}
\item Établir l'équation de la réaction. Quel est le rôle du zinc ?
\item Calculer la masse et la quantité d'ions argent contenus dans le bain.
\item \initiative Quelle masse minimale de zinc faut-il utiliser pour récupérer tout l'argent ?
\item Proposer un test pour vérifier que le bain traité contient des ions zinc(II).
\end{questions}
\end{exercice}
\end{blocLivre}
\begin{blocLivre}{ex-3-15}
\begin{exercice}[3]{Identifier un métal}
\begin{docu}{1}{L'expérience}
Un échantillon de \qty{0,262}{g} d'un métal inconnu est attaqué par de l'acide chlorhydrique en excès. On recueille, par
déplacement d'eau, \qty{96}{mL} de dihydrogène à \qty{20}{\degreeCelsius} et \qty{1013}{hPa}. Le métal est entièrement transformé en ions M$^{2+}$.
\end{docu}
\begin{docu}{2}{Quelques métaux}
\par\smallskip\noindent\hfill\begin{tabular}{lcccc}
\toprule
Métal & magnésium & fer & zinc & cuivre \\
\midrule
$M$ (\unit{g.mol^{-1}}) & 24,3 & 55,8 & 65,4 & 63,5 \\
\bottomrule
\end{tabular}\hfill\null\par\smallskip
Le cuivre n'est pas attaqué par l'acide chlorhydrique.
\end{docu}
\begin{donnees}
$V_\mathrm{m}=\qty{24,0}{L.mol^{-1}}$ ; couple $\ce{H+}/\ce{H2}$.
\end{donnees}
\begin{questions}
\item Établir l'équation de la réaction en notant M le métal.
\item Calculer la quantité de dihydrogène recueillie.
\item \initiative Identifier le métal.
\item Pourquoi le cuivre, de masse molaire proche, peut-il être écarté sans calcul ? Proposer un test qui confirmerait
l'identification.
\end{questions}
\end{exercice}
\end{blocLivre}
